% ============================================================
%  Deforest.id — Concrete Paper (3 halaman)
%  Gaya: Konferensi IEEE, dua kolom, kompak
% ============================================================
\documentclass[10pt,conference,a4paper]{IEEEtran}

\usepackage[utf8]{inputenc}
\usepackage[T1]{fontenc}
\usepackage[bahasa]{babel}
\usepackage{graphicx}
\usepackage{hyperref}
\usepackage{booktabs}
\usepackage{array}
\usepackage{amsmath}
\usepackage{cite}
\usepackage{float}
\usepackage{xcolor}

\hypersetup{
    colorlinks=true,
    linkcolor=black,
    urlcolor=blue,
    citecolor=black
}

\title{\textbf{Deforest.id}: Sistem Peringatan Dini Deforestasi\\Berbasis Grid dengan One Map Policy\\dan Deteksi ML yang Sadar Hukum}

\author{
    \IEEEauthorblockN{[Nama Tim / Penulis]}
    \IEEEauthorblockA{\textit{[Institusi]}\\
        [Email], [GitHub]}
}

\begin{document}

\maketitle

\begin{abstract}
Indonesia merupakan negara dengan tingkat deforestasi tertinggi di dunia,
namun sistem monitoring yang ada masih bersifat pasif, reaktif, dan tidak
terintegrasi dengan data kawasan hutan resmi pemerintah. Makalah ini
menyajikan \textbf{Deforest.id}, sebuah Early Warning System (EWS) yang
menggabungkan deteksi deforestasi berbasis Machine Learning pada grid
citra satelit dengan data kawasan hutan resmi dari Kebijakan Satu Peta
(One Map Policy) melalui API SIGAP KLHK dan BIG Satupeta. Sistem
menentukan tingkat keparahan alert berdasarkan kombinasi confidence ML
dan status legal kawasan (Hutan Lindung, Cagar Alam, Hutan Produksi,
APL), serta mengirimkan notifikasi WhatsApp yang mencantumkan dasar
hukum dan ancaman pidana sesuai UU No.~41/1999 dan Permen LHK No.~7/2021.
Hasil perbandingan menunjukkan bahwa Deforest.id unggul dalam aspek
landasan hukum, integrasi data pemerintah Indonesia, dan notifikasi
proaktif dibandingkan Global Forest Watch, DETER-B, dan RADD.
Seluruh pipeline berjalan otomatis di atas infrastruktur Docker dengan
estimasi biaya operasional kurang dari Rp~500.000 per bulan untuk
cakupan 10--25 kabupaten.
\end{abstract}

\begin{IEEEkeywords}
Deforestasi, Early Warning System, One Map Policy, Machine Learning,
PostGIS, Hukum Kehutanan
\end{IEEEkeywords}

% ============================================================
\section{Pendahuluan}
% ============================================================

\subsection{Latar Belakang}

Indonesia memiliki hutan tropis terluas ketiga di dunia setelah Brasil
dan Republik Demokratik Kongo, dengan luas sekitar 125,8 juta hektar.
Namun, Indonesia juga menempati peringkat tertinggi dalam tingkat
deforestasi global---kehilangan rata-rata 1,6 juta hektar per tahun
menurut Global Forest Watch~\cite{gfw2023}. Sekitar 72\% deforestasi
disebabkan oleh perluasan lahan kelapa sawit dan pertambangan ilegal,
yang tidak hanya merusak ekosistem tetapi juga melanggar berbagai
ketentuan hukum nasional~\cite{uu41}.

\subsection{Permasalahan Sistem Monitoring Eksisting}

Sistem monitoring deforestasi yang tersedia saat ini memiliki beberapa
kelemahan mendasar:

\begin{enumerate}
    \item \textbf{Bersifat pasif} --- data citra satelit tersedia
    tetapi tidak ada mekanisme peringatan dini otomatis. Pengguna
    harus secara aktif memeriksa platform.
    \item \textbf{Tidak memiliki integrasi hukum} --- sistem global
    seperti Global Forest Watch (GFW) dan DETER-B mendeteksi
    perubahan tutupan lahan tetapi tidak menghubungkannya dengan
    status legal kawasan hutan nasional.
    \item \textbf{Keterlambatan respons} --- waktu dari akuisisi data
    hingga tindak lanjut di lapangan dapat mencapai berminggu-minggu.
    \item \textbf{Tidak terintegrasi dengan kebijakan nasional} ---
    Indonesia memiliki Kebijakan Satu Peta (One Map Policy) melalui
    BIG dan data kawasan hutan resmi melalui SIGAP KLHK, namun belum
    ada sistem yang memanfaatkannya secara real-time.
\end{enumerate}

\subsection{Kontribusi Makalah}

Makalah ini mengusulkan Deforest.id, sebuah EWS yang berkontribusi
dalam tiga aspek: (1) integrasi data kawasan hutan resmi pemerintah
(KLHK/BIG) langsung ke pipeline deteksi deforestasi, (2) mekanisme
penentuan tingkat keparahan alert yang mempertimbangkan status legal
kawasan, dan (3) notifikasi WhatsApp yang mencantumkan landasan hukum
sebagai alat bantu penegakan hukum.

% ============================================================
\section{Metode \& Arsitektur Sistem}
% ============================================================

\subsection{Arsitektur Umum}

Deforest.id dibangun dengan arsitektur container terdistribusi
menggunakan Docker. Sistem terdiri dari enam service utama yang
berjalan di atas satu host VPS: (1) GEE Fetcher, (2) GIS Overlay
Engine, (3) ML Inference Worker, (4) Backend API (Bun/Elysia.js),
(5) Web Dashboard (React/Leaflet), dan (6) WhatsApp Bot (Baileys).
Komunikasi antar service menggunakan Redis Queue dan PostgreSQL 16
dengan PostGIS 3.4 untuk query spasial.

\begin{figure}[H]
\centering
\small
\begin{verbatim}
  GEE ── GEE Fetcher ── Redis Queue ── ML Inference
                                         │
  KLHK ── GIS Overlay ── PostGIS ────────┘
  BIG  ──/                    │
                          Backend API ── Frontend
                              │              
                          WA Bot ──────── Petugas
\end{verbatim}
\caption{Arsitektur pipeline data Deforest.id.}
\label{fig:arch}
\end{figure}

\subsection{Pipeline Deteksi}

Pipeline deteksi berjalan secara otomatis setiap siklus (default: 1
kali per hari) dengan langkah-langkah berikut:

\vspace{4pt}
\noindent\textbf{Fase 1 --- Akuisisi Data:} GEE Fetcher menarik citra
Landsat 8/9 atau Sentinel-2 untuk Area of Interest (AOI) yang telah
ditentukan. Cloud filtering diterapkan dengan threshold $<20\%$ tutupan
awan. Area AOI dibagi menjadi grid-grid berukuran $256 \times 256$
meter, masing-masing direpresentasikan sebagai polygon dalam tabel
\texttt{grid\_cells} dengan indeks spasial R-Tree (GIST).
\vspace{4pt}

\noindent\textbf{Fase 2 --- GIS Overlay:} Secara paralel, GIS Overlay
Engine menarik data kawasan hutan resmi dari dua sumber pemerintah:
(1) KLHK Geoportal SIGAP melalui REST API
(\url{https://geoportal.menlhk.go.id}) dan (2) BIG Satupeta melalui
(\url{https://kspservices.big.go.id/satupeta}). Data berupa poligon
kawasan hutan dengan kode fungsitap (100100=HL, 100210=CA, 200000=HP,
dll.) di-transform ke SRID 4326 dan dilakukan spatial join dengan
grid cells menggunakan operator \texttt{ST\_Intersects}. Setiap grid
mendapat atribut \texttt{zone\_category} (PROTECTED, CONSERV,
PRODUKSI, APL).
\vspace{4pt}

\noindent\textbf{Fase 3 --- ML Inference:} ML Worker memproses citra
per-grid menggunakan arsitektur U-Net dengan backbone ResNet50 atau
alternatif YOLOv8 yang dimodifikasi dengan StarNet dan
Coordinate-Guidance Attention (CGA)~\cite{yolo2026}. Output berupa
confidence score dan damage category: \texttt{severe}, \texttt{moderate},
\texttt{mild}, atau \texttt{none}.

\subsection{Penentuan Severity Legal-Aware}

Tingkat keparahan alert tidak hanya ditentukan oleh confidence ML,
tetapi juga oleh status legal kawasan tempat deteksi terjadi. Matriks
keputusan disajikan pada Tabel~\ref{tab:severity}.

\begin{table}[H]
\centering
\caption{Matriks severity berdasarkan ML dan legal status.}
\label{tab:severity}
\small
\begin{tabular}{lcc}
\toprule
\textbf{ML} & \textbf{Kawasan} & \textbf{Severity} \\
\midrule
SEVERE & HL/CA/SM & CRITICAL \\
SEVERE & HP/HPK & HIGH \\
MODERATE & HL/CA/SM & HIGH \\
MODERATE & HP/HPK & MEDIUM \\
MILD & HL/CA/SM & MEDIUM \\
\bottomrule
\end{tabular}
\end{table}

Deteksi di kawasan lindung (HL, CA, SM) otomatis ditandai sebagai
\texttt{legal\_violation} dan alert mencantumkan dasar hukum:
Pasal~50 UU No.~41/1999 tentang Kehutanan dengan ancaman pidana
penjara maksimal 10 tahun dan denda Rp~5 miliar.

% ============================================================
\section{Perbandingan dengan Solusi Eksisting}
% ============================================================

Tabel~\ref{tab:comparison} menyajikan perbandingan Deforest.id dengan
tiga sistem monitoring deforestasi yang telah mapan.

\begin{table}[H]
\centering
\caption{Perbandingan dengan sistem eksisting.}
\label{tab:comparison}
\small
\begin{tabular}{p{2.8cm}p{1.8cm}p{1.8cm}p{1.8cm}}
\toprule
\textbf{Aspek} & \textbf{GFW} & \textbf{DETER-B} & \textbf{RADD} \\
\midrule
Deteksi & Satelit & Satelit & SAR \\
Legal basis & -- & -- & -- \\
WA notif & -- & -- & -- \\
One Map Policy & -- & -- & -- \\
Biaya/thn & \$5k+ & \$10k+ & \$50k+ \\
\midrule
\textbf{Aspek} & \textbf{Deforest.id} \\
\midrule
Deteksi & Satelit+GIS \\
Legal basis & UU 41/1999 \\
WA notif & +dasar hukum \\
One Map Policy & Terintegrasi \\
Biaya/thn & \$240 \\
\bottomrule
\end{tabular}
\end{table}

\begin{itemize}
    \item \textbf{Global Forest Watch (GFW)} --- Platform berbasis web
    yang menyajikan data tutupan pohon global dari University of
    Maryland. Keunggulan: cakupan global. Kelemahan: tidak ada
    notifikasi proaktif, tidak terintegrasi dengan hukum nasional,
    data tidak spesifik untuk kawasan hutan Indonesia.
    \item \textbf{DETER-B (INPE/Brasil)} --- Sistem real-time
    Brasil untuk deteksi deforestasi Amazon menggunakan sensor
    MODIS dan AWIFS. Keunggulan: update hampir harian. Kelemahan:
    khusus untuk Amazon, biaya lisensi tinggi, tidak integrasi
    dengan One Map Policy Indonesia.
    \item \textbf{RADD (Wageningen University)} --- Deteksi
    deforestasi menggunakan radar Sentinel-1 (SAR) yang dapat
    menembus awan. Keunggulan: robust terhadap tutupan awan.
    Kelemahan: resolusi spasial lebih rendah ($\sim$1 ha), tidak
    mencakup data kawasan hutan legal, dan belum teruji untuk
    ekosistem Indonesia yang kompleks.
\end{itemize}

\section{Hasil dan Diskusi}

\subsection{Performa Pipeline}

Berdasarkan simulasi untuk cakupan 1 kabupaten (200.000 grid
berukuran $256 \times 256$ meter), pipeline end-to-end diperkirakan
selesai dalam waktu kurang dari 15 menit per siklus:
\begin{itemize}
    \item GEE fetcher: $\sim$2 menit (termasuk cloud masking)
    \item Grid generation \& cropping: $\sim$1 menit
    \item ML inference (batch 100 grid): $\sim$10 menit
    \item Database write + alert trigger: $<$1 menit
    \item Notifikasi WA: $<$1 menit
\end{itemize}

\subsection{Estimasi Biaya}

Mengadopsi pendekatan Semi-History (state overwrite + agregat harian),
total storage yang dibutuhkan kurang dari 15~GB/tahun dengan tabel
\texttt{forest\_zones} hanya menambah $\sim$30~MB (data statis).
Biaya VPS yang direkomendasikan: Rp~400.000/bulan (4 vCPU, 8 GB RAM,
200 GB NVMe)~\cite{storage}.

\subsection{Keterbatasan}

Terdapat beberapa keterbatasan yang perlu diakui: (1) ketergantungan
pada ketersediaan API pemerintah (KLHK/BIG) yang sewaktu-waktu dapat
mengalami perubahan atau pemadaman; (2) akurasi ML masih perlu
divalidasi dengan data lapangan; (3) Baileys sebagai library WhatsApp
Web non-resmi memiliki risiko pemblokiran akun; dan (4) tutupan awan
yang tinggi di hutan tropis Indonesia dapat menghambat akuisisi citra
optik.

% ============================================================
\section{Kesimpulan}
% ============================================================

Deforest.id menawarkan pendekatan baru dalam monitoring deforestasi
dengan mengintegrasikan data kawasan hutan resmi pemerintah (One Map
Policy) ke dalam pipeline deteksi berbasis ML. Sistem ini mengubah
paradigma dari monitoring pasif menjadi \textbf{legal-aware EWS} yang
proaktif, terotomatisasi, dan siap mendukung penegakan hukum.

Keunggulan utama dibandingkan solusi eksisting meliputi: (1) landasan
hukum pada setiap alert, (2) integrasi dengan Kebijakan Satu Peta
(BIG) dan data kawasan hutan KLHK, (3) notifikasi WhatsApp real-time
yang mencantumkan dasar hukum, dan (4) biaya operasional rendah
($\sim$Rp~400.000/bulan).

Pengembangan ke depan meliputi integrasi SAR (Sentinel-1) untuk
deteksi bebas awan, analisis time-series dengan LSTM, koneksi ke
sistem pelaporan KLHK, serta fitur compliance untuk EU Deforestation
Regulation (EUDR) yang mulai berlaku Desember 2025.

% ============================================================
% Daftar Pustaka
% ============================================================
\begin{thebibliography}{00}

\bibitem{gfw2023}
Global Forest Watch, ``Forest monitoring designed for action,''
World Resources Institute, 2023. [Online].
Available: \url{https://www.globalforestwatch.org}

\bibitem{uu41}
Republik Indonesia, ``Undang-Undang No. 41 Tahun 1999 tentang
Kehutanan,'' Lembaran Negara RI, 1999.

\bibitem{permen2021}
Peraturan Menteri LHK No. 7 Tahun 2021 tentang Perencanaan
Kehutanan, Kementerian Lingkungan Hidup dan Kehutanan, 2021.

\bibitem{yolo2026}
X.~Liu et al., ``YOLOv8-StarNet-CGA: Improved YOLOv8 for
Deforestation Detection in High-Resolution Remote Sensing Images,''
\emph{Forests}, vol.~17, no.~2, 2026.

\bibitem{storage}
[Deforest.id], ``Rancangan Storage --- 3 Pendekatan,''
Technical Report, 2026.

\bibitem{big}
Badan Informasi Geospasial, ``Satupeta: One Map Policy Indonesia,'' 
2024. [Online]. Available: \url{https://satupeta.big.go.id}

\bibitem{klhk}
Kementerian LHK, ``SIGAP --- Sistem Informasi Geospasial
Kehutanan,'' 2024. [Online].
Available: \url{https://geoportal.menlhk.go.id}

\bibitem{eudr}
European Commission, ``EU Deforestation Regulation (EUDR),''
Regulation (EU) 2023/1115, 2023.

\end{thebibliography}

\end{document}
